
\Data\ID{1003138701}
\Updated{2022-04-13 09:04:51}
\Level{A}
\SubArea{Geometric mappings}
\LANGen{
\Question{
When reflecting in a point \( S \), which of lines are mapped to the same lines?}
{Every line passing through the point \( S \).}{Every line that is not passing through the point \( S \).}{There are no lines which are mapped to the same lines in a reflection in the point \( S \).}{Only two mutually perpendicular lines passing through the point \( S \).}{}{}}
\LANGpl{
\Question{
W symetrii środkowej względem punktu  \( S \), które proste są odwzorowywane na te same proste?}
{Każda prosta przechodząca przez punkt \( S \).}{Każda prosta nie przechodząca przez punkt  \( S \).}{Brak prostych odwzorowywanych na te same proste w symetrii środowej względem punktu \( S \).}{Tylko dwie wzajemnie prostopadłe proste przechodzące przez punkt  \( S \).}{}{}}
\LANGcs{
\Question{
Které přímky jsou ve středové souměrnosti dané středem  \( S \) samodružné?}
{Všechny přímky procházející středem souměrnosti  \( S \).}{Všechny přímky, které neprocházejí středem souměrnosti  \( S \).}{Středová souměrnost daná středem \( S \) nemá samodružné přímky.}{Pouze dvě navzájem kolmé přímky procházející středem souměrnosti  \( S \) jsou samodružné.}{}{}}
\LANGsk{
\Question{
Ktoré priamky sú v stredovej súmernosti danej stredom  \( S \) samodružné?}
{Všetky priamky prechádzajúce stredom súmernosti  \( S \).}{Všetky priamky, ktoré neprechádzajú stredom súmernosti  \( S \).}{Stredová súmernosť daná stredom \( S \) nemá samodružné priamky.}{Iba dve navzájom kolmé priamky prechádzajúce stredom súmernosti  \( S \) sú samodružné.}{}{}}
\LANGes{
\Question{
¿Qué rectas se transforman en si mismas por la simetría por el centro  \( S \)?}
{Todas las rectas que pasan por el punto\( S \).}{Todas las rectas que no pasan por el punto \( S \).}{No hay ninguna recta que se transformen en si mismas mediante simetría por el punto \( S \)
.}{Solamente dos rectas perpendiculares que pasan por el punto \( S \).}{}{}}
\End

\Data\ID{1003161101}
\Updated{2021-01-27 03:01:48}
\Level{A}
\SubArea{Geometric mappings}
\LANGen{
\Question{
Given \( p \) is a line of reflection, which lines are mapped to themselves across the line \( p \)?}
{The line \( p \) and all lines perpendicular to \( p \) are mapped to themselves.}{All lines parallel to the line \( p \) are mapped to themselves.}{There are no lines which are mapped to themselves.}{The line \( p \) only.}{}{}}
\LANGpl{
\Question{
Dana jest prosta \( p \). Które proste są odwzorowane są na siebie w symetrii osiowej  względem prostej  \( p \)?}
{Prosta \( p \) oraz wszystkie proste prostopadłe do \( p \) są odwzorowane na siebie.}{Wszystkie proste równoległe do prostej \( p \) są odwzorowane na siebie.}{Brak prostych odwzorowanych na siebie.}{Tylko prosta \( p \).}{}{}}
\LANGcs{
\Question{
Je dána přímka \( p \). Které přímky jsou v osové souměrnosti dané přímkou \( p \) samodružné?}
{Přímky kolmé na osu souměrnosti  \( p \) a přímka \( p \).}{Přímky rovnoběžné s přímkou \( p \).}{Osová souměrnost nemá samodružné přímky.}{Pouze osa souměrnosti  \( p \).}{}{}}
\LANGsk{
\Question{
Je daná priamka \( p \). Ktoré priamky sú v osovej súmernosti danej priamkou \( p \) samodružné?}
{Priamky kolmé na os súmernosti  \( p \) a priamka \( p \).}{Priamky rovnobežné s priamkou \( p \).}{Osová súmernosť nemá samodružné priamky.}{Len os súmernosti  \( p \).}{}{}}
\LANGes{
\Question{
Sea una recta \( p \). ¿Qué rectas se transforman en si mismas por la simetría axial dada por el eje  \( p \) ?}
{Todas las rectas perpendiculares a \( p \) y la recta \( p \).}{Las rectas paralelas a \( p \).}{Ninguna}{Solamente la recta \( p \).}{}{}}
\End

\Data\ID{1103138702}
\Updated{2022-04-13 09:04:34}
\Level{A}
\SubArea{Geometric mappings}
\LANGen{
\Question{
Let \( ABCD \) be a square and let \( M \) be a point, such that \( M\in AB \) and \(|AM|:|MB|=3:1 \). 
Which of the pictures shows the square \( A'B'C'D' \) which is the reflection of the original square \( ABCD \) reflected in the point \( M \)?}
{}{}{}{}{}{}}
\LANGpl{
\Question{
Dany jest kwadrat \( ABCD \) oraz punkt  \( M \),  \( M\in AB \) i \(|AM|:|MB|=3:1 \). 
Na którym z rysunków znajduje się kwadrat  \( A'B'C'D' \), który jest obrazem w symetrii środkowej kwadratu  \( ABCD \) względem punktu  \( M \)?}
{}{}{}{}{}{}}
\LANGcs{
\Question{
Je dán čtverec  \( ABCD \) a bod \( M \), kde \( M \in AB \) a \(|AM|:|MB|=3:1 \). 
Na kterém obrázku je čtverec  \( A'B'C'D' \) obrazem čtverce  \( ABCD \) ve středové souměrnosti se středem v bodě  \( M \)?}
{}{}{}{}{}{}}
\LANGsk{
\Question{
Je daný štvorec  \( ABCD \) a bod \( M \), kde \( M \in AB \) a \(|AM|:|MB|=3:1 \). 
Na ktorom obrázku je štvorec  \( A'B'C'D' \) obrazom štvorca  \( ABCD \) v stredovej súmernosti so stredom v bode  \( M \)?}
{}{}{}{}{}{}}
\LANGes{
\Question{
Sea un cuadrado \( ABCD \) y un punto \( M \), donde  \( M \in AB \) y \(|AM|:|MB|=3:1 \). 
¿En qué dibujo el cuadrado  \( A'B'C'D' \) es la imagen del cuadrado \( ABCD \) por simetría respecto al punto \( M \)?}
{}{}{}{}{}{}}

\IMGanswer{1103138702a_1.pdf}
\IMGanswer{1103138702b_1.pdf}
\IMGanswer{1103138702c.pdf}
\IMGanswer{1103138702d.pdf}\End

\Data\ID{1103138703}
\Updated{2022-04-13 09:04:34}
\Level{A}
\SubArea{Geometric mappings}
\LANGen{
\Question{
Let \( ABCD \)  be a rectangle and let \( M \) be a point, such that \( M\in BD \) and \( |BM|:|MD|=1:3 \). 
Which of the pictures shows the rectangle \( A'B'C'D' \) which is the reflection of the original rectangle \( ABCD \) reflected in the point \( M \)?}
{}{}{}{}{}{}}
\LANGpl{
\Question{
Dany jest prostokąt  \( ABCD \)  oraz punkt \( M \),  \( M\in BD \) i \( |BM|:|MD|=1:3 \). 
Na którym rysunku znajduje się  \( A'B'C'D' \), który jest obrazem w symetrii środkowej prostokąta  \( ABCD \) względem punktu \( M \)?}
{}{}{}{}{}{}}
\LANGcs{
\Question{
Je dán obdélník  \( ABCD \)  a bod  \( M \), kde  \( M \in BD \) a \( |BM|:|MD|=1:3 \). 
Na kterém obrázku je obdélník  \( A'B'C'D' \) obrazem obdélníku \( ABCD \) ve středové souměrnosti se středem v bodě \( M \)?}
{}{}{}{}{}{}}
\LANGsk{
\Question{
Je daný obdĺžnik  \( ABCD \)  a bod  \( M \), kde  \( M \in BD \) a \( |BM|:|MD|=1:3 \). 
Na ktorom obrázku je obdĺžnik  \( A'B'C'D' \) obrazom obdĺžnika \( ABCD \) v stredovej súmernosti so stredom v bode \( M \)?}
{}{}{}{}{}{}}
\LANGes{
\Question{
Sea un rectángulo \( ABCD \) y un punto \( M \), donde  \( M \in BD \) y \( |BM|:|MD|=1:3 \). 
¿En qué dibujo el rectángulo  \( A'B'C'D' \) es la imagen del rectángulo \( ABCD \) mediante simetría respecto al punto \( M \)?}
{}{}{}{}{}{}}

\IMGanswer{1103138703a.pdf}
\IMGanswer{1103138703b.pdf}
\IMGanswer{1103138703c.pdf}
\IMGanswer{1103138703d.pdf}\End

\Data\ID{1103138704}
\Updated{2022-04-13 09:04:32}
\Level{A}
\SubArea{Geometric mappings}
\LANGen{
\Question{
In the given pictures, there are letters \( F \) shown in different positions. Choose the picture that shows a letter \( F \) and its image in a point reflection in the point \( S \).}
{}{}{}{}{}{}}
\LANGpl{
\Question{
Na rysunkach znajduje się litera  \( F \).  Wybierz rysunek na którym znajduje się  litera  \( F \) i jej obraz w symetrii środkowej względem punktu  \( S \).}
{}{}{}{}{}{}}
\LANGcs{
\Question{
Na kterém obrázku je správně zobrazen obraz písmene \( F \) ve středové souměrnosti se středem v bodě \( S \)?}
{}{}{}{}{}{}}
\LANGsk{
\Question{
Na ktorom obrázku je správne zobrazený obraz písmena \( F \) v stredovej súmernosti so stredom v bode \( S \)?}
{}{}{}{}{}{}}
\LANGes{
\Question{
¿En qué dibujo aparece la imagen correcta de la letra \( F \) por simetría respecto al punto \( S \)?}
{}{}{}{}{}{}}

\IMGanswer{1103138704a.pdf}
\IMGanswer{1103138704b.pdf}
\IMGanswer{1103138704c.pdf}
\IMGanswer{1103138704d.pdf}\End

\Data\ID{1103138705}
\Updated{2022-04-13 09:04:01}
\Level{A}
\SubArea{Geometric mappings}
\LANGen{
\Question{
In which of the pictures is the circle \( k’ \) the image of the circle \( k \) in a point reflection in the point \( S \)?}
{}{}{}{}{}{}}
\LANGpl{
\Question{
Na którym rysunku znajduje się okrąg \( k’ \), który jest obrazem okręgu  \( k \) w symetrii środkowej względem punktu  \( S \)?}
{}{}{}{}{}{}}
\LANGcs{
\Question{
Na kterém obrázku je správně zobrazen obraz kružnice \( k \) ve středové souměrnosti se středem v bodě \( S \)?}
{}{}{}{}{}{}}
\LANGsk{
\Question{
Na ktorom obrázku je správne zobrazený obraz kružnice \( k \) v stredovej súmernosti so stredom v bode \( S \)?}
{}{}{}{}{}{}}
\LANGes{
\Question{
¿En qué dibujo aparece la imagen correcta de la circunferencia \( k \) mediante simetría respecto al punto \( S \)?}
{}{}{}{}{}{}}

\IMGanswer{1103138705a.pdf}
\IMGanswer{1103138705b.pdf}
\IMGanswer{1103138705c.pdf}
\IMGanswer{1103138705d.pdf}\End

\Data\ID{1103138706}
\Updated{2022-04-13 09:04:30}
\Level{A}
\SubArea{Geometric mappings}
\LANGen{
\Question{
In which of the pictures is the triangle \( A'B'C' \) the image of  the equilateral triangle \( ABC \) in a point reflection in the point \( S \)?}
{}{}{}{}{}{}}
\LANGpl{
\Question{
Na którym rysunku znajduje się trójkąt \( A'B'C' \), który jest obrazem trójkąta równobocznego \( ABC \) w symetrii środkowej względem punktu  \( S \)?}
{}{}{}{}{}{}}
\LANGcs{
\Question{
Na kterém obrázku je správně zobrazen obraz \( A'B'C' \) trojúhelníku \( ABC \) ve středové souměrnosti se středem v bodě \( S \)?}
{}{}{}{}{}{}}
\LANGsk{
\Question{
Na ktorom obrázku je správne zobrazený obraz \( A'B'C' \) trojuholníka \( ABC \) v stredovej súmernosti so stredom v bode \( S \)?}
{}{}{}{}{}{}}
\LANGes{
\Question{
¿En qué dibujo aparece la imagen correcta \( A'B'C' \) del triángulo \( ABC \) por simetría respecto al punto \( S \)?}
{}{}{}{}{}{}}

\IMGanswer{1103138706a.pdf}
\IMGanswer{1103138706b.pdf}
\IMGanswer{1103138706c.pdf}
\IMGanswer{1103138706d.pdf}\End

\Data\ID{1103138707}
\Updated{2022-04-13 09:04:13}
\Level{A}
\SubArea{Geometric mappings}
\LANGen{
\Question{
Let there be a regular pentagon \( ABCDE \) and a point \( S \). Which of the pictures shows the pentagon \( A'B'C'D'E' \) which is the reflection of the original pentagon \( ABCDE \) reflected in the point \( S \)?}
{}{}{}{}{}{}}
\LANGpl{
\Question{
Dany jest pięciokąt równoboczny \( ABCDE \) oraz punkt  \( S \). Na którym rysunku znajduje się pięciokąt  \( A'B'C'D'E' \), który jest obrazem pięciokąta \( ABCDE \) w symetrii środkowej względem punktu  \( S \)?}
{}{}{}{}{}{}}
\LANGcs{
\Question{
Na kterém obrázku je správně zobrazen obraz  \( A'B'C'D'E' \) pravidelného pětiúhelníku  \( ABCDE \) ve středové souměrnosti se středem v bodě \( S \)?}
{}{}{}{}{}{}}
\LANGsk{
\Question{
Na ktorom obrázku je správne zobrazený obraz  \( A'B'C'D'E' \) pravidelného päťuholníka  \( ABCDE \) v stredovej súmernosti so stredom v bode \( S \)?}
{}{}{}{}{}{}}
\LANGes{
\Question{
¿En qué dibujo aparece la imagen correcta \( A'B'C'D' \) del pentágono regular \( ABCD \)  por simetría respecto al punto \( S \)?}
{}{}{}{}{}{}}

\IMGanswer{1103138707a.pdf}
\IMGanswer{1103138707b.pdf}
\IMGanswer{1103138707c.pdf}
\IMGanswer{1103138707d.pdf}\End

\Data\ID{1103161102}
\Updated{2022-04-13 09:04:29}
\Level{A}
\SubArea{Geometric mappings}
\IMG{1103161102.pdf}
\LANGen{
\Question{
How many lines of symmetry does the parallelogram in the picture have?}
{\( 0 \)}{\( 1 \)}{\( 2 \)}{\( 4 \)}{}{}}
\LANGpl{
\Question{
Ile osi symetrii ma dany równoległobok?}
{\( 0 \)}{\( 1 \)}{\( 2 \)}{\( 4 \)}{}{}}
\LANGcs{
\Question{
Kolik os symetrií má rovnoběžník na obrázku?}
{\( 0 \)}{\( 1 \)}{\( 2 \)}{\( 4 \)}{}{}}
\LANGsk{
\Question{
Koľko osí symetrií má rovnobežník na obrázku?}
{\( 0 \)}{\( 1 \)}{\( 2 \)}{\( 4 \)}{}{}}
\LANGes{
\Question{
¿Cuántos ejes de simetría tiene este paralelogramo?}
{\( 0 \)}{\( 1 \)}{\( 2 \)}{\( 4 \)}{}{}}
\End

\Data\ID{1103161103}
\Updated{2022-04-13 09:04:57}
\Level{A}
\SubArea{Geometric mappings}
\IMG{1103161103.pdf}
\LANGen{
\Question{
How many lines of symmetry does the regular nonagon in the picture have?}
{\( 9 \)}{\( 0 \)}{\( 1 \)}{\( 18 \)}{}{}}
\LANGpl{
\Question{
Ile osi symetrii ma dany dziesięciokąt foremny?}
{\( 9 \)}{\( 0 \)}{\( 1 \)}{\( 18 \)}{}{}}
\LANGcs{
\Question{
Kolik os symetrií má pravidelný devítiúhelník na obrázku?}
{\( 9 \)}{\( 0 \)}{\( 1 \)}{\( 18 \)}{}{}}
\LANGsk{
\Question{
Koľko osí symetrií má pravidelný deväťuholník na obrázku?}
{\( 9 \)}{\( 0 \)}{\( 1 \)}{\( 18 \)}{}{}}
\LANGes{
\Question{
¿Cuántos ejes de simetría tiene este nonágono?}
{\( 9 \)}{\( 0 \)}{\( 1 \)}{\( 18 \)}{}{}}
\End

\Data\ID{1103161104}
\Updated{2022-04-13 09:04:41}
\Level{A}
\SubArea{Geometric mappings}
\IMG{1103161104_0.pdf}
\LANGen{
\Question{
In the picture, there is an equilateral triangle and a line \( o \).  From the next pictures choose the one that shows the triangle and its image in a line reflection across the line \( o \).}
{}{}{}{}{}{}}
\LANGpl{
\Question{
Dany jest trójkąt równoboczny oraz prosta \( o \).  Z podanych rysunków wybierz trójkąt i jego obraz w symetrii osiowej względem prostej  \( o \).}
{}{}{}{}{}{}}
\LANGcs{
\Question{
Je dán rovnostranný trojúhelník a přímka  \( o \) (viz obrázek).  Na kterém obrázku je obraz tohoto trojúhelníku v osové souměrnosti dané přímkou \( o \)?}
{}{}{}{}{}{}}
\LANGsk{
\Question{
Je daný rovnostranný trojuholník a priamka  \( o \) (viď obrázok).  Na ktorom obrázku je obraz tohto trojuholníka v osovej súmernosti danej priamkou \( o \)?}
{}{}{}{}{}{}}
\LANGes{
\Question{
Sea un triángulo equilátero y la recta \( o \) (vea el dibujo). ¿En qué dibujo aparece la imagen del triángulo en simetría respecto a la recta \( o \)?}
{}{}{}{}{}{}}

\IMGanswer{1103161104a.pdf}
\IMGanswer{1103161104b.pdf}
\IMGanswer{1103161104c.pdf}
\IMGanswer{1103161104d.pdf}\End

\Data\ID{1103161105}
\Updated{2022-04-13 09:04:17}
\Level{A}
\SubArea{Geometric mappings}
\IMG{1103161105.pdf}
\LANGen{
\Question{
In the picture, there is letter F and a line \( o \). From the next pictures choose the one that shows an image of F in a line symmetry across the line \( o \).}
{}{}{}{}{}{}}
\LANGpl{
\Question{
Na rysunku znajduje się litera  F oraz prosta  \( o \). Z podanych rysunków wybierz literę  obraz litery F w symetrii osiowej względem prostej \( o \).}
{}{}{}{}{}{}}
\LANGcs{
\Question{
Je dáno písmeno F a přímka \( o \). Na kterém obrázku je obraz písmene F v osové souměrnosti dané přímkou \( o \)?}
{}{}{}{}{}{}}
\LANGsk{
\Question{
Je dané písmeno F a priamka \( o \). Na ktorom obrázku je obraz písmena F v osovej súmernosti danej priamkou \( o \)?}
{}{}{}{}{}{}}
\LANGes{
\Question{
Dada la letra F y la recta \( o \) (observa la imagen). ¿En qué dibujo aparece la imagen de la letra F por simetría respecto a la recta \( o \) ?}
{}{}{}{}{}{}}

\IMGanswer{1103161105a.pdf}
\IMGanswer{1103161105b.pdf}
\IMGanswer{1103161105c.pdf}
\IMGanswer{1103161105d.pdf}\End

\Data\ID{1103161106}
\Updated{2022-05-23 02:05:03}
\Level{A}
\SubArea{Geometric mappings}
\IMG{1103161106.pdf}
\LANGen{
\Question{
Which of the pictures shows a mirror image of the given geometric shape (see the picture) in a mirror reflection across the line \( p \)?}
{}{}{}{}{}{}}
\LANGpl{
\Question{
Na którym rysunku znajduje się obraz podanej figury geometrycznej w symetrii osiowej względem prostej \( p \)?}
{}{}{}{}{}{}}
\LANGcs{
\Question{
Na kterém obrázku je správně zobrazen obrazec v osové souměrnosti dané přímkou \( p \)?}
{}{}{}{}{}{}}
\LANGsk{
\Question{
Na ktorom obrázku je správne zobrazený obrazec v osovej súmernosti danej priamkou \( p \)?}
{}{}{}{}{}{}}
\LANGes{
\Question{
¿En qué dibujo aparece la imágen del primer dibujo respecto al eje de simetría \( p \)?}
{}{}{}{}{}{}}

\IMGanswer{1103161106a.pdf}
\IMGanswer{1103161106b.pdf}
\IMGanswer{1103161106c.pdf}
\IMGanswer{1103161106d.pdf}\End

\Data\ID{1103161107}
\Updated{2022-05-23 02:05:32}
\Level{A}
\SubArea{Geometric mappings}
\IMG{1103161107_0.pdf}
\LANGen{
\Question{
In the picture below there is a trapezoid and a line \( p \). From the next pictures choose the one that shows an image of the trapezoid in a line reflection across the line \( p \).}
{}{}{}{}{}{}}
\LANGpl{
\Question{
Dany jest trapez oraz prosta \( p \).  Z podanych rysunków wybierz  jego obraz w symetrii osiowej względem prostej \( p \).}
{}{}{}{}{}{}}
\LANGcs{
\Question{
Na kterém obrázku je obraz lichoběžníka v osové souměrnosti dané přímkou  \( p \)?}
{}{}{}{}{}{}}
\LANGsk{
\Question{
Na ktorom obrázku je obraz lichobežníka v osovej súmernosti danej priamkou  \( p \)?}
{}{}{}{}{}{}}
\LANGes{
\Question{
¿En qué dibujo aparece  la imagen del trapezoide respecto al eje de simetría  \( p \)?}
{}{}{}{}{}{}}

\IMGanswer{1103161107a.pdf}
\IMGanswer{1103161107b.pdf}
\IMGanswer{1103161107c.pdf}
\IMGanswer{1103161107d.pdf}\End

\Data\ID{1103161108}
\Updated{2021-01-26 01:01:52}
\Level{A}
\SubArea{Geometric mappings}
\LANGen{
\Question{
Which of the traffic signs in the picture does not have any line of symmetry?}
{}{}{}{}{}{}}
\LANGpl{
\Question{
Które ze znaków drogowych nie maja osi symetrii?}
{}{}{}{}{}{}}
\LANGcs{
\Question{
Která dopravní značka nemá osu souměrnosti?}
{}{}{}{}{}{}}
\LANGsk{
\Question{
Ktorá dopravná značka nemá os súmernosti?}
{}{}{}{}{}{}}
\LANGes{
\Question{
¿Cuál de las señales de tráfico no tiene eje de simetría?}
{}{}{}{}{}{}}

\IMGanswer{1103161108b_0.pdf}
\IMGanswer{1103161108a_0.pdf}
\IMGanswer{1103161108c.pdf}
\IMGanswer{1103161108d.pdf}\End

\Data\ID{2010016709}
\Updated{2022-11-02 09:11:38}
\Level{A}
\SubArea{Geometric mappings}
\LANGen{
\Question{
Which of the following answers contains three letters so that each letter is without a point of symmetry? (A letter has a point of symmetry if there exists a point such that the reflection through this point maps the letter into itself.)}
{A, U, K}{S, O, H}{N, U, K}{A, O, H}{}{}}
\LANGpl{
\Question{
Która z poniższych odpowiedzi zawiera trzy litery, takie aby żadna litera nie miała środka symetrii? (Litera ma środek symetrii, jeśli istnieje taki punkt, że obrazem litery względem tego punktu jest ta sama litera.)}
{A, U, K}{S, O, H}{N, U, K}{A, O, H}{}{}}
\LANGcs{
\Question{
Vyberte z odpovědí tu, ve které není ani jedno písmeno středově souměrné. (Písmeno je středově souměrné, pokud existuje středová souměrnost, ve které je písmeno obrazem sebe sama.)}
{A, U, K}{S, O, H}{N, U, K}{A, O, H}{}{}}
\LANGsk{
\Question{
Určte, pre ktorú trojicu písmen nie je ani jedno písmeno stredovo súmerné. (Písmeno je stredovo súmerné, ak existuje stredová súmernosť, v ktorej je písmeno obrazom samého seba.)}
{A, U, K}{S, O, H}{N, U, K}{A, O, H}{}{}}
\LANGes{
\Question{
¿Cuál de las siguientes respuestas contiene tres letras que no presentan una simetría central? (Una letra presenta una simetría central si existe un punto respecto del cual se transforma a si misma.)}
{A, U, K}{S, O, H}{N, U, K}{A, O, H}{}{}}
\End

\Data\ID{2010016710}
\Updated{2022-11-02 09:11:38}
\Level{A}
\SubArea{Geometric mappings}
\LANGen{
\Question{
How many points of symmetry there exist for an equilateral triangle? (A point is the point of symmetry of a triangle if the reflection through this point maps the triangle into itself.)}
{none}{one}{two}{three}{}{}}
\LANGpl{
\Question{
Ile środków symetrii ma trójkąt równoboczny? (Punkt jest środkiem symetrii trójkąta, jeżeli obrazem trójkąta w symetrii względem  tego punktu jest ten sam trójkąt.)}
{Zero}{Jeden}{Dwa}{Trzy}{}{}}
\LANGcs{
\Question{
Kolik středů souměrnosti má rovnostranný trojúhelník? (V takové středové souměrnosti se trojúhelník zobrazí sám na sebe.)}
{žádný}{jeden}{dva}{tři}{}{}}
\LANGsk{
\Question{
Koľko stredov súmernosti má rovnostranný trojuholník? (V takejto stredovej súmernosti sa trojuholník zobrazí sám do seba.)}
{žiadny}{jeden}{dva}{tri}{}{}}
\LANGes{
\Question{
¿Cuántos puntos de simetría existen para un triángulo equilátero? (Un punto A es el punto de simetría de un triángulo si, al calcular el simétrico de cada punto del triángulo respecto de A, se obtiene el mismo triángulo.)}
{ninguno}{uno}{dos}{tres}{}{}}
\End

\Data\ID{9000149301}
\Updated{2022-06-11 07:06:35}
\Level{A}
\SubArea{Geometric mappings}
\LANGen{
\Question{
Which of the following is true for a reflection through a line?}
{It is a congruent mapping (preserves shapes and dimensions).}{It is a similar mapping (preserves shapes, may change dimensions and position).}{It is an identity mapping (objects are mapped to itself).}{}{}{}}
\LANGpl{
\Question{
Wskaż zdanie prawdziwe dotyczące symetrii osiowej.}
{Przekształcenie zachowuje kształty i wymiary.}{Przekształcenie zachowuje kształty, może zmienić wymiary i pozycje.}{Obiekty są odwzorowane na siebie.}{}{}{}}
\LANGcs{
\Question{
Mezi jaká zobrazení patří osová souměrnost?}
{shodná}{podobná}{jde o identitu}{}{}{}}
\LANGsk{
\Question{
Medzi aké zobrazenia patrí osová súmernosť?}
{zhodné}{podobné}{ide o identitu}{}{}{}}
\LANGes{
\Question{
¿Qué tipo de transformación es una simetría central?}
{igualdad}{semejanza}{es la identidad}{}{}{}}
\End

\Data\ID{9000149302}
\Updated{2022-06-11 07:06:57}
\Level{A}
\SubArea{Geometric mappings}
\LANGen{
\Question{
Given a line in a plane, find how many points are mapped onto itself with reflection
through this line.}
{infinitely many}{none}{one}{two}{}{}}
\LANGpl{
\Question{
Dana jest prosta na płaszczyźnie. Wskaż ile punktów jest odwzorowanych na siebie przez symetrię osiową względem tej prostej.}
{nieskończenie wiele}{żaden}{jeden}{dwa}{}{}}
\LANGcs{
\Question{
Kolik samodružných bodů má osová souměrnost?}
{nekonečně mnoho (všechny body osy souměrnosti)}{žádný}{jeden}{právě dva}{}{}}
\LANGsk{
\Question{
Koľko samodružných bodov má osová súmernosť?}
{nekonečne mnoho (všetky body osy súmernosti)}{žiadny}{jeden}{práve dva}{}{}}
\LANGes{
\Question{
¿Cuántos puntos que se transformen en si mismos tiene la simetría axial?}
{infinitos puntos (todos los puntos en el eje de simetría)}{ninguno}{uno}{exactamente dos}{}{}}
\End

\Data\ID{9000149303}
\Updated{2022-06-11 07:06:40}
\Level{A}
\SubArea{Geometric mappings}
\LANGen{
\Question{
Which of the following answers contains three letters with point symmetry? (A letter
has a point of symmetry if there exists a point, such that the reflection through this
point maps the letter into itself.)}
{O, N, Z}{A, M, O}{A, O, B}{H, O, U}{}{}}
\LANGpl{
\Question{
Wskaż, które trzy litery mają środek symetrii?}
{O, N, Z}{A, M, O}{A, O, B}{H, O, U}{}{}}
\LANGcs{
\Question{
Určete, pro kterou trojici písmen platí, že všechna uvedená písmena
mají střed souměrnosti.}
{O, N, Z}{A, M, O}{A, O, B}{H, O, U}{}{}}
\LANGsk{
\Question{
Určte, pre ktorú trojicu písmen platí, že všetky uvedené písmená
majú stred súmernosti.}
{O, N, Z}{A, M, O}{A, O, B}{H, O, U}{}{}}
\LANGes{
\Question{
¿Para qué grupo de letras se cumple que todas las letras tienen un centro de simetría?}
{O, N, Z}{A, M, O}{A, O, B}{H, O, U}{}{}}
\End

\Data\ID{9000149304}
\Updated{2022-04-23 05:04:11}
\Level{A}
\SubArea{Geometric mappings}
\IMG{001493_04-figure0.pdf}
\LANGen{
\Question{
How many points of symmetry there exist for a rhombus? A rhombus -- opposite sides are parallel and all sides have equal length. (A point is a point of
symmetry of the rhombus if the reflection through this point maps the rhombus into
itself.)}
{one}{none}{two}{four}{}{}}
\LANGpl{
\Question{
Ile środków symetrii ma romb?}
{jeden}{nie ma}{dwa}{cztery}{}{}}
\LANGcs{
\Question{
Kolik středů souměrnosti má kosočtverec?}
{jeden}{žádný}{dva}{čtyři}{}{}}
\LANGsk{
\Question{
Koľko stredov súmernosti má kosoštvorec?}
{jeden}{žiadny}{dva}{štyri}{}{}}
\LANGes{
\Question{
¿Cuántos centros de simetría tiene el rombo?}
{uno}{ninguna}{dos}{cuatro}{}{}}
\End

\Data\ID{1103162401}
\Updated{2022-05-23 02:05:10}
\Level{B}
\SubArea{Geometric mappings}
\IMG{1103162401.pdf}
\LANGen{
\Question{
Suppose a square \( ABCD \) rotates about its vertex \( B \). Which of the given pictures shows the square \( A'B'C'D' \) that is the image of the square \( ABCD \) rotated \( 180 \) degrees?}
{}{}{}{}{}{}}
\LANGpl{
\Question{
Kwadrat  \( ABCD \) został obrócony wokół wierzchołka  \( B \). Na którym rysunku znajduje się kwadrat \( A'B'C'D' \), który jest obrazem kwadratu \( ABCD \) w obrocie o  \( 180 \) stopni?}
{}{}{}{}{}{}}
\LANGcs{
\Question{
Na kterém obrázku je správně sestrojen obraz \( A'B'C'D' \) čtverce  \( ABCD \) v otočení daném středem \( B \) a úhlem rotace \( 180 \) stupňů?}
{}{}{}{}{}{}}
\LANGsk{
\Question{
Na ktorom obrázku je správne zostrojený obraz \( A'B'C'D' \) štvorca  \( ABCD \) v otočení danej stredom \( B \) a uhlom rotácie \( 180 \) stupňov?}
{}{}{}{}{}{}}
\LANGes{
\Question{
¿En qué dibujo aparece la imagen correcta \( A'B'C'D' \) del cuadrado \( ABCD \)  en simetría respecto al punto \( S \) y con rotación de \( +180^{\circ} \)?}
{}{}{}{}{}{}}

\IMGanswer{1103162401a.pdf}
\IMGanswer{1103162401b.pdf}
\IMGanswer{1103162401c.pdf}
\IMGanswer{1103162401d.pdf}\End

\Data\ID{1103162402}
\Updated{2022-05-29 12:05:45}
\Level{B}
\SubArea{Geometric mappings}
\IMG{1103162402.pdf}
\LANGen{
\Question{
The rectangle \( ABCD \) rotates \( 90^{\circ} \) anticlockwise about the centre \( S \). Which of the pictures shows the rectangle \( A'B'C'D' \) that is the image of rotated rectangle \( ABCD \)?}
{}{}{}{}{}{}}
\LANGpl{
\Question{
Prostokąt  \( ABCD \) obrócono o  \( 90^{\circ} \) przeciwnie do ruchu wskazówek zegara względem środka \( S \). Na którym rysunku znajduje się prostokąt \( A'B'C'D' \), który jest obrazem  \( ABCD \)?}
{}{}{}{}{}{}}
\LANGcs{
\Question{
Na kterém obrázku je správně sestrojen obraz \( A'B'C'D' \) obdélníku  \( ABCD \) v otočení daném středem \( S \) a úhlem otočení \( +90^{\circ} \)?}
{}{}{}{}{}{}}
\LANGsk{
\Question{
Na ktorom obrázku je správne zostrojený obraz \( A'B'C'D' \) obdĺžnika  \( ABCD \) v otočení danom stredom \( S \) a uhlom otočenia\( +90^{\circ} \)?}
{}{}{}{}{}{}}
\LANGes{
\Question{
¿En qué dibujo aparece la imagen correcta \( A'B'C'D' \) del rectángulo \( ABCD \)  en simetría respecto al punto \( S \) y con rotación de \( +90^{\circ} \)?}
{}{}{}{}{}{}}

\IMGanswer{1103162402a.pdf}
\IMGanswer{1103162402b.pdf}
\IMGanswer{1103162402c.pdf}
\IMGanswer{1103162402d.pdf}\End

\Data\ID{1103162403}
\Updated{2022-05-23 02:05:12}
\Level{B}
\SubArea{Geometric mappings}
\IMG{1103162403.pdf}
\LANGen{
\Question{
An equilateral triangle \( ABC \)  rotates \( 120^{\circ} \) clockwise about the vertex \( A \). Which of the pictures shows the triangle \( A'B'C' \) that is the image of the rotated triangle \( ABC \)?}
{}{}{}{}{}{}}
\LANGpl{
\Question{
Trójkąt równoboczny \( ABC \)  obrócono o \( 120^{\circ} \) zgodnie z  ruchem wskazówek zegara względem wierzchołka  \( A \). Na którym rysunku znajduje się trójkąt \( A'B'C' \), który jest obrazem  trójkąta \( ABC \)?}
{}{}{}{}{}{}}
\LANGcs{
\Question{
Na kterém obrázku je správně sestrojen obraz \( A'B'C' \) rovnostranného trojúhelníku \( ABC \)  v otočení daném středem  \( A \) a úhlem otočení  \( -120^{\circ} \)?}
{}{}{}{}{}{}}
\LANGsk{
\Question{
Na ktorom obrázku je správne zostrojený obraz \( A'B'C' \) rovnostranného trojuholníka \( ABC \)  v otočení danom stredom  \( A \) a uhlom otočenia  \( -120^{\circ} \)?}
{}{}{}{}{}{}}
\LANGes{
\Question{
¿En qué dibujo está la imagen correcta \( A'B'C' \) del triángulo equilátero \( ABC \)  en simetría respecto al punto \( A \) y con rotación de \( -120^{\circ} \)?}
{}{}{}{}{}{}}

\IMGanswer{1103162403a.pdf}
\IMGanswer{1103162403b.pdf}
\IMGanswer{1103162403c.pdf}
\IMGanswer{1103162403d.pdf}\End

\Data\ID{1103162404}
\Updated{2022-05-23 02:05:37}
\Level{B}
\SubArea{Geometric mappings}
\IMG{1103162404_0.pdf}
\LANGen{
\Question{
An equilateral triangle \( ABC \) rotates \( 60^{\circ} \) anticlockwise about its centre \( S \). Which of the pictures shows the triangle \( A'B'C' \) that is the image of the rotated triangle \( ABC \)?}
{}{}{}{}{}{}}
\LANGpl{
\Question{
Trójkąt równoboczny  \( ABC \) obrócono o \( 60^{\circ} \) przeciwnie do ruchu wskazówek zegara względem środka  \( S \). Na którym rysunku znajduj się trójkąt \( A'B'C' \), który jest obrazem  trójkąta \( ABC \)?}
{}{}{}{}{}{}}
\LANGcs{
\Question{
Na kterém obrázku je správně sestrojen obraz \( A'B'C' \) rovnostranného trojúhelníku  \( ABC \) v otočení daném středem \( S \) a úhlem otočení \( +60^{\circ} \)?}
{}{}{}{}{}{}}
\LANGsk{
\Question{
Na ktorom obrázku je správne zostrojený obraz \( A'B'C' \) rovnostranného trojuholníka  \( ABC \) v otočení danom stredom \( S \) a uhlom otočenia \( +60^{\circ} \)?}
{}{}{}{}{}{}}
\LANGes{
\Question{
¿En qué dibujo está la imagen \( A'B'C' \) del triángulo equilátero  \( ABC \)  en giro respecto al punto \( S \) y con el ángulo de giro de \( +60^{\circ} \)?}
{}{}{}{}{}{}}

\IMGanswer{1103162404a_0.pdf}
\IMGanswer{1103162404b_0.pdf}
\IMGanswer{1103162404c_0.pdf}
\IMGanswer{1103162404d_0.pdf}\End

\Data\ID{1103162405}
\Updated{2022-11-01 09:11:41}
\Level{B}
\SubArea{Geometric mappings}
\IMG{1103162405.pdf}
\LANGen{
\Question{
An equilateral triangle \( ABC \) rotates \( 90^{\circ} \) anticlockwise about the point \( O(4;1) \). Which of the pictures shows the triangle \( A'B'C' \) that is the image of the rotated triangle \( ABC \)?}
{}{}{}{}{}{}}
\LANGpl{
\Question{
Trójkąt równoboczny  \( ABC \) obrócono o  \( 90^{\circ} \) przeciwnie do ruchu wskazówek zegara względem punktu  \( O(4;1) \). Na którym rysunku znajduje się trójkąt  \( A'B'C' \), który jest obrazem trójkąta \( ABC \)?}
{}{}{}{}{}{}}
\LANGcs{
\Question{
Na kterém obrázku je správně sestrojen obraz \( A'B'C' \) rovnostranného trojúhelníku \( ABC \) v otočení daném středem v bodě \( O(4;1) \) a úhlem otočení \( +90^{\circ} \)?}
{}{}{}{}{}{}}
\LANGsk{
\Question{
Na ktorom obrázku je správne zostrojený obraz \( A'B'C' \) rovnostranného trojuholníka \( ABC \) v otočení danom stredom v bode \( O(4;1) \) a uhlom otočenia \( +90^{\circ} \)?}
{}{}{}{}{}{}}
\LANGes{
\Question{
¿En qué dibujo aparece la imagen \( A'B'C' \) del triángulo equilátero  \( ABC \)  mediante un giro respecto al punto \( O(4;1) \) y con el ángulo de giro de \( +90^{\circ} \)?}
{}{}{}{}{}{}}

\IMGanswer{1103162405a.pdf}
\IMGanswer{1103162405b.pdf}
\IMGanswer{1103162405c.pdf}
\IMGanswer{1103162405d.pdf}\End

\Data\ID{1103162501}
\Updated{2022-05-23 02:05:46}
\Level{B}
\SubArea{Geometric mappings}
\IMG{1103162501.pdf}
\LANGen{
\Question{
Let a square \( ABCD \) be translated by vector \( \overrightarrow{AC} \). From the pictures below choose the one that shows correctly the translation \( A'B'C'D' \) of the square \( ABCD \).}
{}{}{}{}{}{}}
\LANGpl{
\Question{
Kwadrat \( ABCD \) został przesunięty o wektor \( \overrightarrow{AC} \).  Z podanych rysunków wybierz ten, który pokazuje poprawne przesunięcie \( A'B'C'D' \) kwadratu \( ABCD \).}
{}{}{}{}{}{}}
\LANGcs{
\Question{
Čtverec \( ABCD \) zobrazíme  v posunutí daném vektorem  \( \overrightarrow{AC} \). Vyberte obrázek, na kterém je správně sestrojen obraz \( A'B'C'D' \) čtverce \( ABCD \).}
{}{}{}{}{}{}}
\LANGsk{
\Question{
Štvorec \( ABCD \) zobrazíme  v posunutí daným vektorom  \( \overrightarrow{AC} \). Vyberte obrázok, na ktorom je správne zostrojený obraz \( A'B'C'D' \) štvorca \( ABCD \).}
{}{}{}{}{}{}}
\LANGes{
\Question{
Vamos a trasladar el cuadrado \( ABCD \) mediante el vector \( \overrightarrow{AC} \). Elige el dibujo en el cuál está correctamente la imagen  \( A'B'C'D' \) del cuadrado \( ABCD \).}
{}{}{}{}{}{}}

\IMGanswer{1103162501a_0.pdf}
\IMGanswer{1103162501b_0.pdf}
\IMGanswer{1103162501c_0.pdf}
\IMGanswer{1103162501d_0.pdf}\End

\Data\ID{1103162502}
\Updated{2022-05-23 02:05:19}
\Level{B}
\SubArea{Geometric mappings}
\IMG{1103162502.pdf}
\LANGen{
\Question{
Let a rectangle \( ABCD \) be translated by vector \( \overrightarrow{SB} \), where \( S \) is the midpoint of the rectangle. Which of the following pictures shows correctly the image \( A'B'C'D' \) of the rectangle \( ABCD \) in described translation?}
{}{}{}{}{}{}}
\LANGpl{
\Question{
Prostokąt  \( ABCD \) został przesunięty o wektor \( \overrightarrow{SB} \),  \( S \) to środek prostokąta. Z podanych rysunków wybierz ten, który pokazuje poprawne przesunięcie  \( A'B'C'D' \) prostokąta \( ABCD \)?}
{}{}{}{}{}{}}
\LANGcs{
\Question{
Obdélník \( ABCD \) se středem v bodě  \( S \) zobrazíme v posunutí daném vektorem \( \overrightarrow{SB} \). Vyberte obrázek, na kterém je správně sestrojen obraz \( A'B'C'D' \) obdélníku \( ABCD \).}
{}{}{}{}{}{}}
\LANGsk{
\Question{
Obdĺžnik \( ABCD \) so stredom v bode  \( S \) zobrazíme v posunutí daným vektorom \( \overrightarrow{SB} \). Vyberte obrázok, na ktorom je správne zostrojený obraz \( A'B'C'D' \) obdĺžnika \( ABCD \).}
{}{}{}{}{}{}}
\LANGes{
\Question{
Vamos a trasladar el rectángulo \( ABCD \) con el centro en el punto \( S \) mediante el vector \( \overrightarrow{SB} \). Elige el dibujo donde está la imagen \( A'B'C'D' \) correcta del rectángulo \( ABCD \).}
{}{}{}{}{}{}}

\IMGanswer{1103162502a_0.pdf}
\IMGanswer{1103162502b_0.pdf}
\IMGanswer{1103162502c_0.pdf}
\IMGanswer{1103162502d_0.pdf}\End

\Data\ID{1103162503}
\Updated{2022-05-23 02:05:59}
\Level{B}
\SubArea{Geometric mappings}
\IMG{1103162503.pdf}
\LANGen{
\Question{
A triangle \( ABC \) is translated by vector \( \overrightarrow{CA} \). Which of the following pictures shows correctly the translation \( A'B'C' \) of the triangle?}
{}{}{}{}{}{}}
\LANGpl{
\Question{
Trójkąt \( ABC \) przesunięto o wektor \( \overrightarrow{CA} \). Z podanych rysunków wybierz ten, który pokazuje poprawne przesunięcie  \( A'B'C' \) trójkąta?}
{}{}{}{}{}{}}
\LANGcs{
\Question{
Trojúhelník \( ABC \) zobrazíme v posunutí daném vektorem  \( \overrightarrow{CA} \). Vyberte obrázek, na kterém je správně sestrojen obraz \( A'B'C' \) tohoto trojúhelníku.}
{}{}{}{}{}{}}
\LANGsk{
\Question{
Trojuholník \( ABC \) zobrazíme v posunutí danom vektorom  \( \overrightarrow{CA} \). Vyberte obrázok, na ktorom je správne zostrojený obraz \( A'B'C' \) tohoto trojuholníka.}
{}{}{}{}{}{}}
\LANGes{
\Question{
Vamos a trasladar la imagen de \( ABC \) mediante el vector \( \overrightarrow{CA} \). Elige el dibujo con la imagen \( A'B'C' \) correcta de este triángulo.}
{}{}{}{}{}{}}

\IMGanswer{1103162503a.pdf}
\IMGanswer{1103162503b.pdf}
\IMGanswer{1103162503c.pdf}
\IMGanswer{1103162503d.pdf}\End

\Data\ID{1103162504}
\Updated{2022-05-23 02:05:36}
\Level{B}
\SubArea{Geometric mappings}
\IMG{1103162504.pdf}
\LANGen{
\Question{
Let a triangle \( ABC \) be translated by vector \( \overrightarrow{BT} \), where \( T \) is the centroid of the triangle (see the picture). Which of the following pictures shows correctly the image \( A'B'C' \) of the triangle \( ABC \) in described?}
{}{}{}{}{}{}}
\LANGpl{
\Question{
Trójkąt \( ABC \) przesunięto o wektor \( \overrightarrow{BT} \),  \( T \) to punkt przecięcia środkowych trójkąta (spójrz na rysunek). Który z podanych rysunków poprawnie wskazuje przesunięcie  \( A'B'C' \) trójkąta \( ABC \)?}
{}{}{}{}{}{}}
\LANGcs{
\Question{
Je dán trojúhelník  \( ABC \) a jeho těžiště  \( T \) (viz obrázek). Na kterém obrázku je správně sestrojen obraz  \( A'B'C' \) tohoto trojúhelníku v posunutí daném vektorem \( \overrightarrow{BT} \)?}
{}{}{}{}{}{}}
\LANGsk{
\Question{
Je daný trojuholník  \( ABC \) a jeho ťažisko  \( T \) (viď obrázok). Na ktorom obrázku je správne zostrojený obraz  \( A'B'C' \) tohoto trojuholníka v posunutí danom vektorom \( \overrightarrow{BT} \)?}
{}{}{}{}{}{}}
\LANGes{
\Question{
Sea un triángulo  \( ABC \) y su baricentro \( T \) (vea el dibujo). ¿En qué dibujo aparece la imagen \( A'B'C' \) correcta de este triángulo hecha por trasladarlo mediante el vector \( \overrightarrow{BT} \)?}
{}{}{}{}{}{}}

\IMGanswer{1103162504a.pdf}
\IMGanswer{1103162504b.pdf}
\IMGanswer{1103162504c.pdf}
\IMGanswer{1103162504d.pdf}\End

\Data\ID{1103162505}
\Updated{2022-05-23 02:05:29}
\Level{B}
\SubArea{Geometric mappings}
\IMG{1103162505.pdf}
\LANGen{
\Question{
The triangle \( ABC \) is translated by the vector \( \vec{u} \) (see the picture). Which of the following pictures shows correctly the image \( A'B'C' \) of the triangle \( ABC \) in the given translation?}
{}{}{}{}{}{}}
\LANGpl{
\Question{
Trójkąt \( ABC \) został przesunięty o wektor \( \vec{u} \) (spójrz na rysunek). Który z podanych rysunków poprawnie wskazuje przesunięcie \( A'B'C' \) trójkąta \( ABC \)?}
{}{}{}{}{}{}}
\LANGcs{
\Question{
Trojúhelník \( ABC \) zobrazíme v posunutí daném vektorem  \( \vec{u} \) (viz obrázek). Vyberte obrázek, na kterém je správně sestrojen obraz  \( A'B'C' \) trojúhelníku \( ABC \) v tomto posunutí?}
{}{}{}{}{}{}}
\LANGsk{
\Question{
Trojuholník \( ABC \) zobrazíme v posunutí danom vektorom  \( \vec{u} \) (viď obrázok). Vyberte obrázok, na ktorom je správne zostrojený obraz  \( A'B'C' \) trojuholníka \( ABC \) v tomto posunutí?}
{}{}{}{}{}{}}
\LANGes{
\Question{
Vamos a trasladar \( ABC \) mediante el vector  \( \vec{u} \) (observa el dibujo). Elige el dibujo con la imagen \( A'B'C' \) correcta de este triángulo.}
{}{}{}{}{}{}}

\IMGanswer{1103162505a.pdf}
\IMGanswer{1103162505b.pdf}
\IMGanswer{1103162505c.pdf}
\IMGanswer{1103162505d.pdf}\End

\Data\ID{2110016701}
\Updated{2022-11-02 09:11:00}
\Level{B}
\SubArea{Geometric mappings}
\IMG{klon-figure0.pdf}
\LANGen{
\Question{
Let \(ABCD\) be a rectangle where \(|AB|=2|BC|\) and let \(S\) be a point on the perpendicular bisector of the side \(BC\) such that its distance from \(BC\) is \(\frac12|BC|\) (see the picture). The rectangle \(ABCD\) rotates \(90^{\circ}\) clockwise about \(S\). Which of the pictures shows the rectangle \(A'B'C'D'\) that is the image of rotated rectangle \(ABCD\)?}
{}{}{}{}{}{}}
\LANGpl{
\Question{
Niech \(ABCD\) będzie prostokątem, gdzie \(|AB|=2|BC|\) i niech \(S\) będzie punktem na dwusiecznej prostopadłej boku \(BC\) takim, że jego odległość od \(BC\) to \(\frac12|BC|\) (patrz rysunek). Prostokąt \(ABCD\) obraca się \(90^{\circ}\) zgodnie z ruchem wskazówek zegara o \(S\). Który z rysunków przedstawia prostokąt \(A'B'C'D'\), który jest obrazem obróconego prostokąta \(ABCD\)?}
{}{}{}{}{}{}}
\LANGcs{
\Question{
Je dán obdélník \(ABCD\), \(|AB|=2|BC|\). Dále je dán bod \(S\), který leží na ose strany \(BC\), jeho vzdálenost od \(BC\) je \(\frac12|BC|\) (viz obrázek). Obdélník \(ABCD\) otočíme o \(90^{\circ}\) v záporném smyslu kolem \(S\). Který obdélník \(A'B'C'D'\) zobrazuje takto otočený obdélník \(ABCD\)?}
{}{}{}{}{}{}}
\LANGsk{
\Question{
Daný je obdĺžnik \(ABCD\), \(|AB|=2|BC|\) a bod \(S\), ktorý leží na osi strany \(BC\) a jeho vzdialenosť od \(BC\) je \(\frac12|BC|\) (pozri obrázok). Obdĺžnik \(ABCD\) otočíme o \(90^{\circ}\) v zápornom zmysle okolo bodu \(S\). Ktorý obdĺžnik \(A'B'C'D'\) zobrazuje takto otočený obdĺžnik \(ABCD\)?}
{}{}{}{}{}{}}
\LANGes{
\Question{
Sea \(ABCD\) un rectángulo donde \(|AB|=2|BC|\) y sea \(S\) un punto en la mediatriz del lado \(BC\) tal que su distancia a \(BC\) es \(\frac12|BC|\) (ver la imagen). El rectángulo \(ABCD\) gira \(90^{\circ}\) en el sentido horario alrededor de \(S\). ¿Cuál de las siguientes imágenes muestra el rectángulo \(A'B'C'D'\) que es la imagen del rectángulo girado \(ABCD\)?}
{}{}{}{}{}{}}

\IMGanswer{klon-figure1_0.pdf}
\IMGanswer{klon-figure2_0.pdf}
\IMGanswer{klon-figure3_0.pdf}
\IMGanswer{klon-figure4.pdf}\End

\Data\ID{2110016702}
\Updated{2022-11-02 09:11:00}
\Level{B}
\SubArea{Geometric mappings}
\IMG{klon-figure5_0.pdf}
\LANGen{
\Question{
An equilateral triangle \( ABC \) rotates \( 90^{\circ} \) clockwise about the point \( O(4;2) \). Which of the pictures shows the triangle \( A'B'C' \) that is the image of the rotated triangle \( ABC \)?}
{}{}{}{}{}{}}
\LANGpl{
\Question{
Trójkąt równoboczny \( ABC \) obraca się \( 90^{\circ} \) zgodnie z ruchem wskazówek zegara wokół punktu \( O(4;2) \). Który z obrazków przedstawia trójkąt \( A'B'C' \), który jest obrazem obróconego trójkąta \( ABC \)?}
{}{}{}{}{}{}}
\LANGcs{
\Question{
Rovnostranný trojúhelník \( ABC \) otočíme o \( 90^{\circ} \)  v záporném smyslu kolem bodu \( O(4;2) \). Který obrázek ukazuje trojúhelník \( A'B'C' \), který je obrazem trojúhelníka \( ABC \)?}
{}{}{}{}{}{}}
\LANGsk{
\Question{
Rovnostranný trojuholník \( ABC \) otočíme o \( 90^{\circ} \)  v zápornom zmysle okolo bodu \( O(4;2) \). Ktorý trojúholník \( A'B'C' \) zobrazuje takto otočený trojuholník \( ABC \)?}
{}{}{}{}{}{}}
\LANGes{
\Question{
Un triángulo equilátero \( ABC \) gira \( 90^{\circ} \) en sentido horario alrededor del punto \( O(4;2) \). ¿Cuál de las siguientes imágenes muestra el triángulo \( A'B'C' \) que corresponde a la imagen del triángulo girado \( ABC \)?}
{}{}{}{}{}{}}

\IMGanswer{klon-figure6_0.pdf}
\IMGanswer{klon-figure7_0.pdf}
\IMGanswer{klon-figure8_0.pdf}
\IMGanswer{klon-figure9_0.pdf}\End

\Data\ID{2110016703}
\Updated{2022-11-02 09:11:03}
\Level{B}
\SubArea{Geometric mappings}
\IMG{klon-figure10.pdf}
\LANGen{
\Question{
Let a rectangle \( ABCD \) be translated by vector \( \overrightarrow{AS} \), where \( S \) is the center of the rectangle. Which of the following pictures shows correctly the image \( A'B'C'D' \) of the rectangle \( ABCD \) in described translation?}
{}{}{}{}{}{}}
\LANGpl{
\Question{
Niech prostokąt \( ABCD \) zostanie przesunięty o wektor \( \overrightarrow{AS} \), gdzie \( S \) jest środkiem prostokąta. Który z poniższych obrazków pokazuje poprawnie obraz \( A'B'C'D' \) prostokąta \( ABCD \) w opisanym przesunięciu?}
{}{}{}{}{}{}}
\LANGcs{
\Question{
Obdélník \( ABCD \) posuneme o vektor \( \overrightarrow{AS} \), kde \( S \) je střed obdélníku. Který z obrázků znázorňuje obraz \( A'B'C'D' \) obdélníku \( ABCD \) v daném posunutí?}
{}{}{}{}{}{}}
\LANGsk{
\Question{
Obdĺžnik \( ABCD \) posunieme o vektor \( \overrightarrow{AS} \), kde \( S \) je stred daného obdĺžnika. Ktorý z obrázkov znázorňuje obdĺžnik \( A'B'C'D' \) ako obraz obdĺžnika \( ABCD \) v danom posunutí?}
{}{}{}{}{}{}}
\LANGes{
\Question{
Sea el rectángulo \( ABCD \) trasladado por el vector \( \overrightarrow{AS} \), donde \( S \) es el centro del rectángulo. ¿Cuál de las siguientes imágenes muestra correctamente el rectángulo \( A'B'C'D' \) que corresponde a la traslación descrita del rectángulo \( ABCD \)?}
{}{}{}{}{}{}}

\IMGanswer{klon-figure11.pdf}
\IMGanswer{klon-figure12.pdf}
\IMGanswer{klon-figure13.pdf}
\IMGanswer{klon-figure14.pdf}\End

\Data\ID{2110016704}
\Updated{2022-11-02 09:11:03}
\Level{B}
\SubArea{Geometric mappings}
\IMG{klon-figure15.pdf}
\LANGen{
\Question{
The triangle \( ABC \) is translated by the vector \( \vec{u} \) (see the picture). Which of the following pictures shows correctly the image \( A'B'C' \) of the triangle \( ABC \) in the given translation?}
{}{}{}{}{}{}}
\LANGpl{
\Question{
Trójkąt \( ABC \) jest przesunięty o wektor \( \vec{u} \) (patrz rysunek). Który z poniższych rysunków poprawnie przedstawia obraz \( A'B'C' \) trójkąta \( ABC \) w podanym przesunięciu?}
{}{}{}{}{}{}}
\LANGcs{
\Question{
Trojúhelník \( ABC \) posuneme o vektor \( \vec{u} \) (viz obrázek). Který z obrázků znázorňuje správně obraz \( A'B'C' \) trojúhelníku \( ABC \) v daném posunutí?}
{}{}{}{}{}{}}
\LANGsk{
\Question{
Trojuholník \( ABC \) posunieme o vektor \( \vec{u} \) (pozri obrázok). Ktorý z obrázkov znázorňuje trojuholník \( A'B'C' \) ako obraz trojuholníka \( ABC \) v danom posunutí?}
{}{}{}{}{}{}}
\LANGes{
\Question{
El triángulo \( ABC \) es trasladado por el vector \( \vec{u} \) (ver la imagen). ¿Cuál de las siguientes imágenes muestra correctamente el triángulo \( A'B'C' \) que corresponde a la traslación descrita del triángulo \( ABC \)?}
{}{}{}{}{}{}}

\IMGanswer{klon-figure16.pdf}
\IMGanswer{klon-figure17.pdf}
\IMGanswer{klon-figure18.pdf}
\IMGanswer{klon-figure19.pdf}\End

\Data\ID{9000149305}
\Updated{2021-01-27 02:01:29}
\Level{B}
\SubArea{Geometric mappings}
\LANGen{
\Question{
Given a translation \(T\) 
of a plane, find the lines which are mapped to the same line by
\(T\).}
{All lines parallel to the translation vector are mapped into itself.}{All lines perpendicular to the translation vector are mapped into itself.}{There are no lines which are mapped into itself by the translation.}{Every line is mapped into itself by the translation.}{}{}}
\LANGpl{
\Question{
Dane jest przesunięcie  \(T\) 
na płaszczyźnie, wskaż proste, które są odwzorowane na te same proste za pomocą
\(T\).}
{Wszystkie proste równoległe do wektora przesunięcia są odwzorowane na siebie.}{Wszystkie proste prostopadłe do wektora przesunięcia są odwzorowane na siebie.}{Brak prostych odwzorowanych na siebie za pomocą przesunięcia.}{Każda prosta jest odwzorowywana na siebie za pomocą przesunięcia.}{}{}}
\LANGcs{
\Question{
Které přímky jsou v posunutí samodružné?}
{Všechny přímky, které jsou rovnoběžné se směrem posunutí.}{Osa souměrnosti posunutí.}{Všechny přímky, které jsou kolmé na směr posunutí.}{Posunutí nemá samodružné přímky.}{}{}}
\LANGsk{
\Question{
Ktoré priamky sú v posunutí samodružné?}
{Všetky priamky, ktoré sú rovnobežné so smerom posunutia.}{Os súmernosti posunutia.}{Všetky priamky, ktoré sú kolmé na smer posunutia.}{Posunutie nemá samodružné priamky.}{}{}}
\LANGes{
\Question{
¿Qué rectas se transforman a si mismas por traslación?}
{Todas las rectas paralelas a la dirección de la traslación.}{El eje de la simetría de la traslación}{Todas las rectas perpendiculares a la dirección de la traslación.}{Ninguna recta}{}{}}
\End

\Data\ID{9000149306}
\Updated{2022-06-11 07:06:16}
\Level{B}
\SubArea{Geometric mappings}
\LANGen{
\Question{
Given a translation of a plane, find the property of a line obtained by translating a line
\(r\). The
line \(r\) is
neither parallel not perpendicular to the translation vector.}
{The resulting line is parallel to the line
\(r\).}{The resulting line is perpendicular to the translation vector.}{The resulting line is perpendicular to the line
\(r\).}{The resulting line is the line \(r\).
(The line \(r\) 
is mapped into itself.)}{}{}}
\LANGpl{
\Question{
Obrazem prostej  \(r\) w przesunięciu o wektor, który nie jest ani równoległy, ani prostopadły do prostej  \(r\) jest:}
{Prosta jest równoległa do prostej
\(r\).}{Prosta jest prostopadła do wektora przesunięcia.}{Prosta jest prostopadła do prostej
\(r\).}{Prosta jest prostą \(r\).
(Prosta \(r\) 
jest odwzorowana na siebie.)}{}{}}
\LANGcs{
\Question{
Obrazem přímky \(r\),
která není rovnoběžná se směrem posunutí, ani kolmá na směr
posunutí, je:}
{přímka rovnoběžná s danou přímkou
\(r\)}{přímka kolmá na směr posunutí}{přímka kolmá k dané přímce
\(r\)}{samodružná přímka}{}{}}
\LANGsk{
\Question{
Obrazom priamky \(r\),
ktorá nie je rovnobežná so smerom posunutia, ani kolmá na smer
posunutia, je:}
{priamka rovnobežná s danou priamkou
\(r\)}{priamka kolmá na smer posunutia}{priamka kolmá k danej priamke
\(r\)}{samodružná priamka}{}{}}
\LANGes{
\Question{
Sea una recta \(r\) que no es paralela ni perpendicular a la dirección de traslación. Su imagen es:}
{una recta paralela a \(r\)}{una recta perpendicular a la dirección de traslación}{una recta perpendicular a \(r\)}{la misma recta \(r\)}{}{}}
\End

\Data\ID{9000149307}
\Updated{2022-06-11 07:06:46}
\Level{B}
\SubArea{Geometric mappings}
\LANGen{
\Question{
The rotation through an angle \(\alpha = 180^{\circ }\) 
is equivalent to some other geometric mapping. Which one?}
{reflection through the point}{reflection through the line}{translation}{}{}{}}
\LANGpl{
\Question{
Obrót o kąt  \(\alpha = 180^{\circ }\) 
jest odpowiednikiem innego przekształcenia geometrycznego. Wskaż, którego?}
{symetrii środkowej}{symetrii osiowej}{przesunięciu}{}{}{}}
\LANGcs{
\Question{
Otočení o úhel \(\alpha = 180^{\circ }\) je ekvivalentní jinému zobrazení. Určete, kterému.}
{středové souměrnosti}{osové souměrnosti}{posunutí}{}{}{}}
\LANGsk{
\Question{
Otočenie o uhol \(\alpha = 180^{\circ }\) je ekvivalentné inému zobrazeniu. Určite, ktorému.}
{stredovú súmernosť}{osovú súmernosť}{posunutie}{}{}{}}
\LANGes{
\Question{
Un giro de ángulo \(\alpha = 180^{\circ }\) es equivalente a otra transformación. ¿Cuál?}
{simetría central}{simetría axial}{traslación}{}{}{}}
\End

\Data\ID{9000149308}
\Updated{2022-06-11 07:06:05}
\Level{B}
\SubArea{Geometric mappings}
\LANGen{
\Question{
Consider a rotation through an angle either
\(\alpha = 180^{\circ }\) or
\(\alpha = 360^{\circ }\). How
many lines which are mapped into itself exists for this rotation?}
{infinitely many}{none}{one}{two}{}{}}
\LANGpl{
\Question{
Rozważmy obrót o kąt
\(\alpha = 180^{\circ }\) lub
\(\alpha = 360^{\circ }\). Ile prostych jest odwzorowanych na siebie?}
{nieskończenie wiele}{żadna}{jedna}{dwie}{}{}}
\LANGcs{
\Question{
Kolik přímek je v otočení samodružných, je-li velikost úhlu otočení
\(\alpha = 180^{\circ }\) nebo
\(\alpha = 360^{\circ }\)?}
{nekonečně mnoho (všechny přímky procházející středem otočení)}{žádná}{právě jedna (přímka procházející středem otočení)}{právě dvě}{}{}}
\LANGsk{
\Question{
Koľko priamok je v otočení samodružných, ak veľkosť uhla otočenia
\(\alpha = 180^{\circ }\) nebo
\(\alpha = 360^{\circ }\)?}
{nekonečne veľa (všetky priamky prechádzajúce stredom otočenia)}{žiadne}{práve jedna (priamka prechádzajúca stredom otočenia)}{práve dve}{}{}}
\LANGes{
\Question{
¿Cuántas rectas se transforman en si mismas por rotación si los ángulos de rotación son
\(\alpha = 180^{\circ }\) o
\(\alpha = 360^{\circ }\)?}
{infinitas (todas las rectas que pasan por el centro de rotación)}{ninguna}{exactamente una (la que pasa por el centro de rotación)}{exactamente dos}{}{}}
\End

\Data\ID{1103162601}
\Updated{2022-05-23 02:05:11}
\Level{C}
\SubArea{Geometric mappings}
\IMG{1103162601.pdf}
\LANGen{
\Question{
Let \( ABCD \) be a square. Find the dilated image of the square with \( A \) being the centre of dilation and the scale factor \( \frac12 \)  (see the picture).}
{}{}{}{}{}{}}
\LANGpl{
\Question{
Dany jest kwadrat \( ABCD \). Wskaż obraz kwadratu uzyskanego po dylatacji, gdzie \( A \) to środek dylatacji,  a współczynnik skali to  \( \frac12 \)  (spójrz na rysunek).}
{}{}{}{}{}{}}
\LANGcs{
\Question{
Je dán čtverec \( ABCD \). Na kterém obrázku je obraz čtverce \( ABCD \) ve stejnolehlosti se středem v bodě  \( A \) a koeficientem \( \frac12 \)?}
{}{}{}{}{}{}}
\LANGsk{
\Question{
Je daný štvorec \( ABCD \). Na ktorom obrázku je obraz štvorca \( ABCD \) v rovnoľahlosti so stredom v bode  \( A \) a koeficientom \( \frac12 \)?}
{}{}{}{}{}{}}
\LANGes{
\Question{
Sea un cuadrado \( ABCD \). ¿En qué dibujo aparece la imagen correcta de  \( ABCD \) hecha por homotecia con el centro en el punto \( A \) y el coeficiente \( \frac12 \)?}
{}{}{}{}{}{}}

\IMGanswer{1103162601a.pdf}
\IMGanswer{1103162601b.pdf}
\IMGanswer{1103162601c.pdf}
\IMGanswer{1103162601d.pdf}\End

\Data\ID{1103162602}
\Updated{2022-05-29 12:05:45}
\Level{C}
\SubArea{Geometric mappings}
\IMG{1103162602.pdf}
\LANGen{
\Question{
Let \( ABCD \) be a square. Find the dilated image of the square with \( S \) being the centre of dilation and scale factor \( -\frac12 \). The point \( S \) is also the middle of the square \( ABCD \) (see the picture).}
{}{}{}{}{}{}}
\LANGpl{
\Question{
Dany jest kwadrat \( ABCD \). Wskaż obraz kwadratu uzyskanego po dylatacji, gdzie \( S \) to środek dylatacji a współczynnik skali jest równy \( -\frac12 \). Punkt \( S \) to również środek kwadratu \( ABCD \) (spójrz na rysunek).}
{}{}{}{}{}{}}
\LANGcs{
\Question{
Je dána čtverec  \( ABCD \). Na kterém obrázku je obraz čtverce  \( ABCD \) ve stejnolehlosti se středem v bodě  \( S \) a koeficientem \( -\frac12 \), kde bod \( S \) je středem čtverce \( ABCD \)?}
{}{}{}{}{}{}}
\LANGsk{
\Question{
Je daný štvorec  \( ABCD \). Na ktorom obrázku je obraz štvorca  \( ABCD \) v rovnoľahlosti so stredom v bode  \( S \) a koeficientom \( -\frac12 \), kde bod \( S \) je stredom štvorca \( ABCD \)?}
{}{}{}{}{}{}}
\LANGes{
\Question{
Sea un cuadrado \( ABCD \). ¿En quél dibujo aparece la imagen correcta de  \( ABCD \) hecha por homotecia con el centro en el punto \( S \) y el coeficiente \( -\frac12 \), donde el punto \( S \) es el centro del cuadro?}
{}{}{}{}{}{}}

\IMGanswer{1103162602a.pdf}
\IMGanswer{1103162602b.pdf}
\IMGanswer{1103162602c.pdf}
\IMGanswer{1103162602d.pdf}\End

\Data\ID{1103162603}
\Updated{2022-05-29 12:05:45}
\Level{C}
\SubArea{Geometric mappings}
\IMG{1103162603.pdf}
\LANGen{
\Question{
Let \( ABC \) be a triangle (see the picture). Find the dilated image of the triangle with \( A \) being the centre of dilation and the scale factor \( \frac54 \).}
{}{}{}{}{}{}}
\LANGpl{
\Question{
Dany jest trójkąt  \( ABC \) (spójrz na rysunek). Wskaż obraz trójkąta uzyskanego po dylatacji, gdzie \( A \) to środek dylatacji a współczynnik skali to  \( \frac54 \).}
{}{}{}{}{}{}}
\LANGcs{
\Question{
Je dán trojúhelník  \( ABC \). Na kterém obrázku je obraz trojúhelníku  \( ABC \) ve stejnolehlosti se středem v bodě  \( A \) a koeficientem  \( \frac54 \)?}
{}{}{}{}{}{}}
\LANGsk{
\Question{
Je daný trojuholník  \( ABC \). Na ktorom obrázku je obraz trojuholníka  \( ABC \) v rovnoľahlosti so stredom v bode  \( A \) a koeficientom  \( \frac54 \)?}
{}{}{}{}{}{}}
\LANGes{
\Question{
Sea un triángulo \( ABC \). ¿En cuál dibujo aparece la imagen correcta de  \( ABC \) hecha por homotecia con el centro en el punto \( A \) y el coeficiente \( \frac54 \)}
{}{}{}{}{}{}}

\IMGanswer{1103162603a_0.pdf}
\IMGanswer{1103162603b_0.pdf}
\IMGanswer{1103162603c_0.pdf}
\IMGanswer{1103162603d_0.pdf}\End

\Data\ID{1103162604}
\Updated{2022-05-29 12:05:45}
\Level{C}
\SubArea{Geometric mappings}
\IMG{1103162604.pdf}
\LANGen{
\Question{
Let \( ABC\) be a triangle with the centroid \( T \) (see the picture). Find the dilated image of the triangle with \( T \) being the centre of dilation and the scale factor \( -\frac12 \).}
{}{}{}{}{}{}}
\LANGpl{
\Question{
Dany jest trójkąt \( ABC \),  \( T \) to  punkt przecięcia środkowych trójkąta (spójrz na rysunek). Wskaż obraz trójkąta uzyskanego po dylatacji, gdzie \( T \) to środek dylatacji, a współczynnik skali to  \( -\frac12 \).}
{}{}{}{}{}{}}
\LANGcs{
\Question{
Je dán trojúhelník  \( ABC \) a jeho těžiště \( T \). Na kterém obrázku je obraz trojúhelníku \( ABC \) ve stejnolehlosti se středem v bodě \( T \) a koeficientem  \( -\frac12 \)?}
{}{}{}{}{}{}}
\LANGsk{
\Question{
Je daný trojuholník  \( ABC \) a jeho ťažisko \( T \). Na ktorom obrázku je obraz trojuholníka \( ABC \) v rovnoľahlosti so stredom v bode \( T \) a koeficientom  \( -\frac12 \)?}
{}{}{}{}{}{}}
\LANGes{
\Question{
Sea un triángulo \( ABC \) y su baricentro \(T \). ¿En qué dibujo aparece la imagen correcta de  \( ABC \) hecha por homotecia con el centro en el punto \( T \) y el coeficiente \( -\frac12 \)?}
{}{}{}{}{}{}}

\IMGanswer{1103162604a_0.pdf}
\IMGanswer{1103162604b_0.pdf}
\IMGanswer{1103162604c_0.pdf}
\IMGanswer{1103162604d_0.pdf}\End

\Data\ID{1103162605}
\Updated{2022-05-29 12:05:45}
\Level{C}
\SubArea{Geometric mappings}
\IMG{1103162605.pdf}
\LANGen{
\Question{
Let \(ABC\) be a triangle (see the picture). Find the dilated image of the triangle with \( O \) (the origin) being the centre of dilation and the scale factor \( 2 \).}
{}{}{}{}{}{}}
\LANGpl{
\Question{
Dany jest trójkąt \(ABC \) (spójrz na rysunek). Wskaż obraz trójkąta  uzyskanego po dylatacji, gdzie \( O \) (środek układu współrzędnych) to środek dylatacji, a współczynnik skali to \( 2 \).}
{}{}{}{}{}{}}
\LANGcs{
\Question{
V souřadném systému je dán trojúhelník \(ABC \). Na kterém obrázku je obraz trojúhelníku \(ABC \) ve stejnolehlosti se středem v bodě  \( O \) (počátku souřadného systému) a koeficientem \( 2 \)?}
{}{}{}{}{}{}}
\LANGsk{
\Question{
V súradnicovom systéme je daný trojuholník \(ABC \). Na ktorom obrázku je obraz trojuholníka \(ABC \) v rovnoľahlosti so stredom v bode  \( O \) (začiatku súradnicového systému) a koeficientom \( 2 \)?}
{}{}{}{}{}{}}
\LANGes{
\Question{
Sea un triángulo \(ABC \) en los ejes de coordenadas. ¿En cuál de los dibujos está la imagen de \(ABC \) por homotecia con el centro en \( O \) (el origen de las coordenadas) y razón \( 2 \)?}
{}{}{}{}{}{}}

\IMGanswer{1103162605a.pdf}
\IMGanswer{1103162605b.pdf}
\IMGanswer{1103162605c.pdf}
\IMGanswer{1103162605d.pdf}\End

\Data\ID{1103162606}
\Updated{2023-12-28 12:12:20}
\Level{C}
\SubArea{Geometric mappings}
\IMG{1103162606.pdf}
\LANGen{
\Question{
The figure below shows two semicircles \( \widehat{AB} \) and \( \widehat{DA} \).  Find the dilation scale factor if you know that the semicircle \( \widehat{DA} \) is the image of the semicircle \( \widehat{AB} \)  in dilation with the centre \( C \).}
{\( -0.5 \)}{\( 0.5 \)}{\( -0.25 \)}{\( 0.75 \)}{}{}}
\LANGpl{
\Question{
Rysunek poniżej przedstawia dwa półkola \( \widehat{AB} \) i \( \widehat{DA} \). Wyznacz wskaźnik dylatacji jeśli wiesz, ze półkole \( \widehat{DA} \) to obraz półkola \( \widehat{AB} \)  ze środkiem dylatacji  \( C \).}
{\( -0{,}5 \)}{\( 0{,}5 \)}{\( -0{,}25 \)}{\( 0{,}75 \)}{}{}}
\LANGcs{
\Question{
Jsou dány dvě půlkružnice  \( \widehat{AB} \) a \( \widehat{DA} \) (viz obrázek).  Určete koeficient stejnolehlosti se středem v bodě \( C \), která zobrazuje půlkružnici \( \widehat{AB} \) na půlkružnici \( \widehat{DA} \).}
{\( -0{,}5 \)}{\( 0{,}5 \)}{\( -0{,}25 \)}{\( 0{,}75 \)}{}{}}
\LANGsk{
\Question{
Na obrázku sú znázornené 2 polkružnice \( \widehat{AB} \) a \( \widehat{DA} \).  Určte koeficient rovnoľahlosti, ak viete, že polkružnica \( \widehat{DA} \) je obrazom polkružnice \( \widehat{AB} \)  v rovnoľahlosti so stredom \( C \).}
{\( -0{,}5 \)}{\( 0{,}5 \)}{\( -0{,}25 \)}{\( 0{,}75 \)}{}{}}
\LANGes{
\Question{
La siguiente imagen representa dos semicircunferencias \( \widehat{AB} \) y \( \widehat{DA} \). Halla la razón de homotecia con el centro \( C \), que representa la semicircunferencia \( \widehat{DA} \) en la semicircunferencia \( \widehat{AB} \).}
{\( -0.5 \)}{\( 0.5 \)}{\( -0.25 \)}{\( 0.75 \)}{}{}}
\End

\Data\ID{2110016705}
\Updated{2022-11-02 09:11:03}
\Level{C}
\SubArea{Geometric mappings}
\IMG{klon-figure20.pdf}
\LANGen{
\Question{
Let \( ABCD \) be a square. Find the dilated image of the square with \( S \) being the centre of dilation and scale factor \( \frac12 \). The point \( S \) is also the middle of the square \( ABCD \) (see the picture).}
{}{}{}{}{}{}}
\LANGpl{
\Question{
Niech \( ABCD \) będzie kwadratem. Znajdź obraz kwadratu względem punktu \( S \) będącym środkiem jednokładności i skali \( \frac12 \). Punkt \( S \) jest jednocześnie środkiem kwadratu \( ABCD \) (patrz rysunek).}
{}{}{}{}{}{}}
\LANGcs{
\Question{
Je dán čtverec \( ABCD \), bod \( S \) je jeho střed. Najděte obraz čtverce ve stejnolehlosti se středem \( S \) a koeficientem \( \frac12 \) (viz obrázek).}
{}{}{}{}{}{}}
\LANGsk{
\Question{
Daný je štvorec \( ABCD \) kde bod \( S \) je jeho stred. Nájdite obraz štvorca v rovnoľahlosti so stredom \( S \) a koeficientom \( \frac12 \) (pozri obrázok).}
{}{}{}{}{}{}}
\LANGes{
\Question{
Sea \( ABCD \) un cuadrado. Halla la imagen dilatada del cuadrado siendo \( S \) el centro de la dilatación y el factor de escala \( \frac12 \). El punto \( S \) corresponde también al centro del cuadrado \( ABCD \) (ver la imagen).}
{}{}{}{}{}{}}

\IMGanswer{klon-figure21.pdf}
\IMGanswer{klon-figure22.pdf}
\IMGanswer{klon-figure23.pdf}
\IMGanswer{klon-figure24.pdf}\End

\Data\ID{2110016706}
\Updated{2022-11-02 09:11:03}
\Level{C}
\SubArea{Geometric mappings}
\IMG{klon-figure25.pdf}
\LANGen{
\Question{
Let \( ABC \) be a triangle (see the picture). Find the dilated image of the triangle with \( B \) being the centre of dilation and the scale factor \( \frac32 \).}
{}{}{}{}{}{}}
\LANGpl{
\Question{
Niech \( ABC \) będzie trójkątem (patrz rysunek). Znajdź obraz trójkąta w jednokładności o środku w punkcie \( B \)  i skali jednokładności \( \frac32 \).}
{}{}{}{}{}{}}
\LANGcs{
\Question{
Najděte obraz trojúhelníku \( ABC \) ve stejnolehlosti se středem v bodě \( B \) a koeficientem \( \frac32 \).}
{}{}{}{}{}{}}
\LANGsk{
\Question{
Najdite obraz trojuholníka \( ABC \) v rovnoľahlosti so stredom v bode \( B \) a koeficientom \( \frac32 \).}
{}{}{}{}{}{}}
\LANGes{
\Question{
Sea \( ABC \) un triángulo (ver la imagen). Halla la figura dilatada del triángulo siendo \( B \) el centro de la dilatación y el factor de escala \( \frac32 \).}
{}{}{}{}{}{}}

\IMGanswer{klon-figure26_0.pdf}
\IMGanswer{klon-figure27.pdf}
\IMGanswer{klon-figure28.pdf}
\IMGanswer{klon-figure29.pdf}\End

\Data\ID{2110016707}
\Updated{2022-11-02 09:11:03}
\Level{C}
\SubArea{Geometric mappings}
\IMG{klon-figure30.pdf}
\LANGen{
\Question{
Let \( ABC\) be a triangle with the centroid \( T \) (see the picture). Find the dilated image of the triangle with \( T \) being the centre of dilation and the scale factor \( \frac12 \).}
{}{}{}{}{}{}}
\LANGpl{
\Question{
Niech \( ABC\) będzie trójkątem o środku ciężkości  \( T \) (patrz rysunek). Znajdź obraz trójkąta w jednokładności o środku w punkcie \( T \) i skali jednokładności \( \frac12 \).}
{}{}{}{}{}{}}
\LANGcs{
\Question{
Určete obraz trojúhelníka \( ABC\) ve stejnolehlosti se středem v těžišti \( T \) trojúhelníka (viz obrázek) a koeficientem stejnolehlosti \( \frac12 \).}
{}{}{}{}{}{}}
\LANGsk{
\Question{
Určte obraz trojuholníka \( ABC\) v rovnoľahlosti so stredom v ťažisku trojuholníka \( T \) (pozri obrázok) a koeficientom rovnoľahlosti \( \frac12 \).}
{}{}{}{}{}{}}
\LANGes{
\Question{
Sea \( ABC\) un triángulo con baricentro \( T \) (ver la imagen). Halla la figura dilatada del triángulo siendo \( T \) el centro de la dilatación y el factor de escala \( \frac12 \).}
{}{}{}{}{}{}}

\IMGanswer{klon-figure31_0.pdf}
\IMGanswer{klon-figure32.pdf}
\IMGanswer{klon-figure33.pdf}
\IMGanswer{klon-figure34.pdf}\End

\Data\ID{2110016708}
\Updated{2022-11-02 09:11:03}
\Level{C}
\SubArea{Geometric mappings}
\IMG{klon-figure35_0.pdf}
\LANGen{
\Question{
Let \(ABC\) be a triangle (see the picture). Find the dilated image of the triangle with \( O \) (the origin) being the centre of dilation and the scale factor \( -2 \).}
{}{}{}{}{}{}}
\LANGpl{
\Question{
Niech \(ABC\) będzie trójkątem (patrz rysunek). Znajdź obraz trójkąta w jednokładności o środku w punkcie  \( O \) (początek układu współrzędnych) i skali jednokładności \( -2 \).}
{}{}{}{}{}{}}
\LANGcs{
\Question{
Určete obraz trojúhelníka \(ABC\) ve stejnolehlosti se středem v počátku soustavy souřadnic \( O \) a koeficientem stejnolehlosti \( -2 \).}
{}{}{}{}{}{}}
\LANGsk{
\Question{
Určte obraz trojuholníka \(ABC\) v rovnoľahlosti so stredom v počiatku súradnicovej sústavy \( O \) a koeficientom rovnoľahlosti \( -2 \).}
{}{}{}{}{}{}}
\LANGes{
\Question{
Sea \(ABC\) un triángulo (ver la imagen). Halla la figura dilatada del triángulo siendo \( O \) (el origen) el centro de la dilatación y el factor de escala \( -2 \).}
{}{}{}{}{}{}}

\IMGanswer{klon-figure36_0.pdf}
\IMGanswer{klon-figure37_1.pdf}
\IMGanswer{klon-figure38_2.pdf}
\IMGanswer{klon-figure39_3.pdf}\End

\Data\ID{9000149309}
\Updated{2022-06-11 07:06:03}
\Level{C}
\SubArea{Geometric mappings}
\LANGen{
\Question{
Consider a dilatation which maps \(A\) 
onto \(B\). The center is
the dilatation is \(S\).
Find a correct statement.}
{The point \(S\) is on the
line through the points \(A\) 
and \(B\).}{The points \(S\),
\(A\) and
\(B\) form a right
triangle \(ABS\).}{The distance from \(S\) 
to \(A\) is smaller than
the distance from \(A\) 
to \(B\).}{The points \(S\),
\(A\) and
\(B\) form a
triangle \(ABS\) 
with at least two sides of equal length.}{}{}}
\LANGpl{
\Question{
Rozważmy izometrię, które odwzorowuje \(A\) 
na  \(B\). Środek izometrii to 
 \(S\).
Wskaż zdanie prawdziwe.}
{Punkt  \(S\) leży na prostej przechodzącej przez punkty
 \(A\) 
i \(B\).}{Punkty \(S\),
\(A\) i
\(B\) tworzą trójkąt prostokątny
 \(ABS\).}{Odległość od  \(S\) 
do \(A\) jest mniejsza niż odległość 
z \(A\) 
do \(B\).}{Punkty \(S\),
\(A\) i
\(B\) tworzą trójkąt \(ABS\), którego przynajmniej dwa boki są równej długości.}{}{}}
\LANGcs{
\Question{
Uvažujme stejnolehlost se středem \(S\), která zobrazuje bod \(A\) na bod
\(B\). Vyberte správné tvrzení.}
{Bod \(S\) leží na přímce \(AB\).}{Body \(A\), \(B\) a \(S\) tvoří pravoúhlý trojúhelník.}{Vzdálenost mezi body \(S\) a \(A\) je menší než vzdálenost mezi body \(S\) a
\(B\).}{Body \(A\), \(B\), a \(S\) tvoří trojúhelník, jehož alespoň dvě strany jsou stejně
dlouhé.}{}{}}
\LANGsk{
\Question{
Uvažujme o rovnoľahlosti so stredom \(S\), ktorá zobrazuje bod \(A\) na bod
\(B\). Vyberte správne tvrdenie.}
{Bod \(S\) leží na priamke \(AB\).}{Body \(A\), \(B\) a \(S\) tvorí pravouhlý trojuholník.}{Vzdialenosť medzi bodmi \(S\) a \(A\) je menšia než vzdialenosť medzi bodmi \(S\) a
\(B\).}{Body \(A\), \(B\), a \(S\) tvoria trojuholník, ktorého aspoň dve strany sú rovnako
dlhé.}{}{}}
\LANGes{
\Question{
Sea una homotecia con centro en  \(S\), que transforma el punto \(A\) en 
\(B\). Elige la declaración correcta.}
{El punto \(S\) pertenece a la recta \(AB\).}{Los puntos \(A\), \(B\) y \(S\) forman un triángulo rectángulo.}{La distancia entre los puntos \(S\) y \(A\) es menor que la distancia entre los puntos \(S\) y
\(B\).}{Los puntos \(A\), \(B\), y \(S\) forman un triángulo dónde como mínimo dos lados tienen la misma longitud.}{}{}}
\End

\Data\ID{9000149310}
\Updated{2022-06-11 07:06:58}
\Level{C}
\SubArea{Geometric mappings}
\LANGen{
\Question{
The dilatation defined by a center and a scale factor
\(k = -1\) is
equivalent to another geometric mapping. Which one?}
{reflection through the point}{translation}{reflection through the line}{rotation}{}{}}
\LANGpl{
\Question{
Wskaż rodzaj symetrii, jeżeli współczynnik jest równy
\(k = -1\).}
{symetria środkowa}{przesunięcie}{symetria osiowa}{obrót}{}{}}
\LANGcs{
\Question{
O jakou souměrnost se jedná, je-li koeficient stejnolehlosti
\(k = -1\)?}
{středovou souměrnost}{posunutí}{osovou souměrnost}{otočení}{}{}}
\LANGsk{
\Question{
O akú súmernosť ide, ak je koeficient rovnoľahlosti
\(k = -1\)?}
{stredovú súmernosť}{posunutie}{osovú súmernosť}{otočenie}{}{}}
\LANGes{
\Question{
¿Qué tipo de simetría tenemos si su razón de semejanza es
\(k = -1\)?}
{simetría por un punto}{trasladación}{simetría por un eje}{rotación}{}{}}
\End
